% Compiler avec : xelatex main.tex (2 passes pour la table des matières)
\documentclass[12pt,a4paper]{article}

% === Polices et langues (XeLaTeX) ===
\usepackage{fontspec}
\setmainfont{Linux Libertine O}   % Police Unicode avec cyrillique
\setsansfont{Linux Libertine O}
\setmonofont[Scale=0.85]{Menlo}    % Monospace avec cyrillique (macOS)

\usepackage{polyglossia}
\setdefaultlanguage{french}
\setotherlanguage{ukrainian}

% Commande pour insérer du texte ukrainien inline
\newcommand{\ukr}[1]{\texturkrainian{#1}}

% === Mise en page ===
\usepackage[margin=2.5cm]{geometry}
\usepackage{parskip}             % Espacement entre paragraphes, pas d'indentation

% === Typographie et couleurs ===
\usepackage{xcolor}
\definecolor{codeblue}{HTML}{2a7ae2}
\definecolor{codegray}{HTML}{555555}
\definecolor{codebg}{HTML}{f5f5f5}
\definecolor{accent-crimson}{HTML}{DC143C}

% === Code source ===
\usepackage{listings}
\lstset{
  basicstyle=\ttfamily\small,
  backgroundcolor=\color{codebg},
  frame=single,
  framerule=0.4pt,
  rulecolor=\color{codegray!30},
  breaklines=true,
  breakatwhitespace=true,
  tabsize=2,
  showstringspaces=false,
  keywordstyle=\color{codeblue}\bfseries,
  commentstyle=\color{codegray}\itshape,
  stringstyle=\color{accent-crimson},
  numbers=left,
  numberstyle=\tiny\color{codegray},
  numbersep=8pt,
  xleftmargin=1.5em,
  framexleftmargin=1.5em,
}
\lstdefinelanguage{svelte}{
  morekeywords={import, from, export, let, const, function, if, else, each, as, await, then, catch},
  morecomment=[l]{//},
  morecomment=[s]{/*}{*/},
  morestring=[b]",
  morestring=[b]',
  sensitive=true,
}
\lstdefinelanguage{javascript}{
  morekeywords={import, from, export, let, const, function, return, if, else, for, switch, case, break, default, new, this, true, false, null, typeof, class, extends},
  morecomment=[l]{//},
  morecomment=[s]{/*}{*/},
  morestring=[b]",
  morestring=[b]',
  morestring=[b]`,
  sensitive=true,
}

% === Images ===
\usepackage{graphicx}
\graphicspath{{img/}}
% Mode draft pour les images manquantes (décommenter si besoin)
% \usepackage[draft]{graphicx}

% === Tableaux ===
\usepackage{booktabs}
\usepackage{array}
\usepackage{longtable}

% === Divers ===
\usepackage{hyperref}
\hypersetup{
  colorlinks=true,
  linkcolor=codeblue,
  urlcolor=codeblue,
  citecolor=codeblue,
  pdfauthor={},
  pdftitle={Documentation — Application vocabulaire ukrainien},
}
\usepackage{enumitem}
\usepackage{fancyvrb}            % Verbatim amélioré
\usepackage{tcolorbox}           % Encadrés colorés
\tcbuselibrary{listings}

% === Encadré pour les notes ===
\newtcolorbox{notebox}{
  colback=codebg,
  colframe=codegray!40,
  fonttitle=\bfseries,
  title=Note,
  sharp corners,
  boxrule=0.4pt,
}

% === Encadré pour les exemples ukrainiens ===
\newtcolorbox{ukrbox}{
  colback=blue!3,
  colframe=codeblue!40,
  sharp corners,
  boxrule=0.4pt,
}

% === Métadonnées ===
\title{%
  \textbf{Documentation technique}\\[0.3em]
  \Large Application web — Vocabulaire ukrainien\\[0.5em]
  \large SvelteKit · Svelte 5 · Vite
}
\author{}
\date{\today}

% ============================================================
\begin{document}

\maketitle
\tableofcontents
\newpage

\section{Présentation générale}

\subsection{Objectif}

Cette application web permet l'apprentissage du vocabulaire ukrainien à travers une interface interactive. Elle présente des mots classés par catégorie grammaticale (noms, verbes, adjectifs, etc.) avec leurs déclinaisons ou conjugaisons complètes, des accents toniques visuels, et des phrases d'exemple.

L'application est construite avec \textbf{SvelteKit} (framework web basé sur Svelte~5) et déployée en site statique via \texttt{adapter-static}.

\subsection{Fonctionnalités principales}

\begin{enumerate}
  \item \textbf{Navigation par catégorie} --- La sidebar gauche liste tous les mots, regroupés par catégorie grammaticale (\ukr{nom}, \ukr{adj.}, \ukr{verbe}, etc.).

  \item \textbf{Fiches détaillées} --- Un clic sur un mot affiche ses formes fléchies dans le panneau central : tableau de déclinaison (noms, adjectifs) ou de conjugaison (verbes).

  \item \textbf{Accents toniques} --- Une case à cocher (en bas à droite) active ou désactive l'affichage des accents. Par exemple, \ukr{оди́н} avec accent vs \ukr{один} sans accent.

  \item \textbf{Info-bulle au survol} --- Le survol d'un mot ukrainien (span \texttt{.ukr}) affiche une bulle avec le lemme accentué et les métadonnées grammaticales (catégorie, cas, nombre, genre).

  \item \textbf{Tableau de grammaire épinglable} --- Un clic sur un mot ukrainien «~épingle~» la sidebar droite, qui affiche le tableau complet de déclinaison/conjugaison. La forme survolée est mise en gras.

  \item \textbf{Couleur par cas grammatical} --- Au survol, la couleur du mot reflète son cas : nominatif (hérité), génitif (vert), accusatif (bleu), locatif (brun), instrumental/datif (rose), vocatif (rouge foncé).

  \item \textbf{Phrases d'exemple} --- Chaque mot peut avoir des phrases associées, affichées sous la fiche détaillée avec traduction française.

  \item \textbf{Recherche de phrases} --- Une page dédiée (\texttt{/phrases}) permet de filtrer les phrases par métadonnées (auteur, cours, source).

  \item \textbf{Génération de conjugaisons} --- Dans les phrases d'exemple, une case à cocher déploie toutes les formes conjuguées d'un verbe pour un temps donné.

  \item \textbf{Responsive} --- L'interface s'adapte aux écrans étroits (empilement des colonnes sous 768px).
\end{enumerate}

\subsection{Captures d'écran}

% ============ PLACEHOLDER ============
% Capture : page principale complète avec la sidebar gauche (liste des mots)
% et le panneau central (détails d'un nom, par exemple "будинок")
% Taille recommandée : pleine largeur du navigateur, ~1200x800px
\begin{figure}[ht]
  \centering
  \fbox{\parbox{0.8\textwidth}{\centering\vspace{3cm}\textit{[screenshot-main.png]\\Page principale : sidebar + détails}\vspace{3cm}}}
  % \includegraphics[width=0.9\textwidth]{screenshot-main.png}
  \caption{Page principale --- liste des mots (gauche) et fiche détaillée (centre).}
  \label{fig:main}
\end{figure}

% ============ PLACEHOLDER ============
% Capture : page /phrases avec la barre de recherche en haut à droite
% et quelques phrases affichées avec leur traduction
\begin{figure}[ht]
  \centering
  \fbox{\parbox{0.8\textwidth}{\centering\vspace{3cm}\textit{[screenshot-phrases.png]\\Page de recherche de phrases}\vspace{3cm}}}
  % \includegraphics[width=0.9\textwidth]{screenshot-phrases.png}
  \caption{Page de recherche de phrases avec filtre par métadonnées.}
  \label{fig:phrases}
\end{figure}

\subsection{Installation et lancement}

\begin{lstlisting}[language=bash, title=Commandes de base]
npm install          # Installer les dependances
npm run dev          # Serveur de developpement (hot reload)
npm run build        # Generer le site statique dans build/
npm run preview      # Previsualiser le site genere
npm run test         # Lancer les tests unitaires (Vitest)
\end{lstlisting}

\newpage
\section{Architecture du projet}

\subsection{Arborescence}

\begin{lstlisting}[title=Structure des fichiers]
appli_web_cc/
|-- src/
|   |-- app.css                    # Styles globaux
|   |-- app.html                   # Template HTML de base
|   |-- routes/
|   |   |-- +layout.svelte         # Layout global (charge les stores)
|   |   |-- +layout.server.js      # Chargement des donnees (prerender)
|   |   |-- +page.svelte           # Page principale (vocabulaire)
|   |   |-- phrases/
|   |       |-- +page.svelte       # Page de recherche de phrases
|   |-- lib/
|       |-- components/            # 13 composants Svelte
|       |   |-- WordList.svelte
|       |   |-- WordDetails.svelte
|       |   |-- NounDetails.svelte
|       |   |-- AdjectiveDetails.svelte
|       |   |-- VerbDetails.svelte
|       |   |-- BaseDetails.svelte
|       |   |-- HtmlContent.svelte
|       |   |-- GrammarSidebar.svelte
|       |   |-- GrammarTable.svelte
|       |   |-- UkrSpan.svelte
|       |   |-- ExamplePhrases.svelte
|       |   |-- PhraseList.svelte
|       |   |-- AccentCheckbox.svelte
|       |-- stores/                # Stores Svelte (etat global)
|       |   |-- dataStore.js       # wordData, phraseData
|       |   |-- uiStore.js         # selectedWord, accentEnabled, ...
|       |-- utils/                 # Fonctions utilitaires
|           |-- accent.js
|           |-- parsing.js
|           |-- dataAccess.js
|           |-- gramFunc.js
|           |-- colors.js
|           |-- i18n.js
|           |-- phrases.js
|           |-- index.js           # Re-export central
|-- static/
|   |-- data.json                  # Dictionnaire complet (genere par Python)
|   |-- phrases.json               # Phrases d'exemple
|-- outil_python/                  # Pipeline de generation des donnees
|   |-- build_entries_2_(nooj_opt).py
|   |-- phrases_a_traiter.json
|   |-- out.json
|   |-- entries_report.html
|-- tests/
|   |-- utils/                     # Tests unitaires (Vitest)
|       |-- accent.test.js
|       |-- parsing.test.js
|       |-- dataAccess.test.js
|       |-- gramFunc.test.js
|       |-- i18n.test.js
|       |-- phrases.test.js
|-- package.json
|-- svelte.config.js               # adapter-static
|-- vite.config.js                 # Vitest + alias $lib
\end{lstlisting}

\subsection{Flux de données}

Le chargement des données suit ce chemin~:

\begin{enumerate}
  \item \texttt{+layout.server.js} lit \texttt{static/data.json} et \texttt{static/phrases.json} au moment du \emph{prerender} (build statique).
  \item Les données sont passées au \texttt{+layout.svelte} via la prop \texttt{data}.
  \item Un \texttt{\$effect} dans le layout initialise les stores globaux \texttt{wordData} et \texttt{phraseData}.
  \item Tous les composants accèdent aux données via ces stores.
\end{enumerate}

\begin{lstlisting}[language=javascript, title=+layout.server.js --- chargement des donnees]
import { readFileSync } from 'fs';
import { resolve } from 'path';

export const prerender = true;

export function load() {
  const wordData = JSON.parse(
    readFileSync(resolve('static/data.json'), 'utf-8')
  );
  const phraseData = JSON.parse(
    readFileSync(resolve('static/phrases.json'), 'utf-8')
  );
  return { wordData, phraseData };
}
\end{lstlisting}

\begin{lstlisting}[language=svelte, title=+layout.svelte --- initialisation des stores]
<script>
  import { wordData, phraseData } from '$lib/stores/dataStore.js';
  let { data, children } = $props();

  $effect(() => {
    wordData.set(data.wordData);
    phraseData.set(data.phraseData);
  });
</script>
{@render children()}
<GrammarSidebar />
<AccentCheckbox />
\end{lstlisting}

\subsection{Graphe de dépendances des composants}

\begin{lstlisting}[title=Arbre des composants]
+layout.svelte
|-- GrammarSidebar
|   |-- GrammarTable
|-- AccentCheckbox
|-- {children}
    |
    |-- +page.svelte (/)
    |   |-- WordList
    |   |   |-- HtmlContent (disableHover=true)
    |   |-- WordDetails
    |       |-- NounDetails
    |       |-- AdjectiveDetails
    |       |-- VerbDetails
    |       |-- BaseDetails
    |       |-- ExamplePhrases
    |           |-- HtmlContent
    |
    |-- phrases/+page.svelte
        |-- PhraseList
            |-- HtmlContent
\end{lstlisting}

Les composants \texttt{GrammarSidebar} et \texttt{AccentCheckbox} sont rendus dans le layout global --- ils sont donc présents sur toutes les pages.

\subsection{Technologies}

\begin{tabular}{ll}
  \toprule
  \textbf{Élément} & \textbf{Technologie} \\
  \midrule
  Framework & SvelteKit 2.0 (Svelte 5.0) \\
  Bundler & Vite 5.0 \\
  Déploiement & \texttt{adapter-static} (site pré-rendu) \\
  Tests & Vitest 1.6.1 + jsdom \\
  Données & JSON statique généré par Python \\
  \bottomrule
\end{tabular}

\newpage
\section{Composants Svelte}

Cette section décrit chacun des 13 composants de l'application, regroupés par domaine fonctionnel. Pour chaque composant, on indique ses \emph{props} (données reçues du parent), les \emph{stores} utilisés, et le rendu produit.

% ----------------------------------------------------------------
\subsection{Navigation et contrôles globaux}
% ----------------------------------------------------------------

\subsubsection{WordList.svelte}

\begin{description}
  \item[Rôle] Sidebar gauche : liste tous les mots du dictionnaire, regroupés par catégorie grammaticale.
  \item[Props] Aucune.
  \item[Stores] \texttt{wordData} (lecture), \texttt{selectedWord} et \texttt{selectedCategory} (écriture au clic).
  \item[Rendu] Pour chaque catégorie non vide, un titre \texttt{<h2>} suivi d'une liste de boutons. Chaque bouton contient un \texttt{HtmlContent} avec \texttt{disableHover=true} --- les mots de la sidebar ne déclenchent pas les info-bulles.
  \item[Interaction] Un clic sur un mot met à jour les stores \texttt{selectedWord} et \texttt{selectedCategory}, ce qui déclenche l'affichage des détails dans le panneau central.
\end{description}

Les catégories sont affichées dans cet ordre fixe~:

\begin{center}
\begin{tabular}{ll}
  \toprule
  \textbf{Clé} & \textbf{Label affiché} \\
  \midrule
  \texttt{nom} & Noms \\
  \texttt{adj} & Adjectifs \\
  \texttt{pron} & Pronoms (déterminants) \\
  \texttt{verb} & Verbes \\
  \texttt{conj} & Conjonctions \\
  \texttt{part} & Particules \\
  \texttt{card} & Cardinaux \\
  \texttt{proposs} & Pronoms possessifs \\
  \texttt{proper} & Pronoms personnels \\
  \texttt{prep} & Prépositions \\
  \bottomrule
\end{tabular}
\end{center}

\subsubsection{AccentCheckbox.svelte}

\begin{description}
  \item[Rôle] Case à cocher flottante (position fixe, en bas à droite) pour activer/désactiver l'affichage des accents toniques.
  \item[Props] Aucune.
  \item[Stores] \texttt{accentEnabled} (lecture et écriture).
  \item[Rendu] Un \texttt{<input type="checkbox">} lié au store. Coché par défaut.
\end{description}

% ----------------------------------------------------------------
\subsection{Fiches détaillées}
% ----------------------------------------------------------------

\subsubsection{WordDetails.svelte}

\begin{description}
  \item[Rôle] Panneau central : affiche la fiche détaillée du mot sélectionné.
  \item[Props] Aucune.
  \item[Stores] \texttt{wordData}, \texttt{selectedWord}, \texttt{selectedCategory} (lecture).
  \item[Rendu] Utilise \texttt{\$derived} pour calculer \texttt{details} à partir des stores. Selon la catégorie, délègue l'affichage à un sous-composant~:
\end{description}

\begin{center}
\begin{tabular}{ll}
  \toprule
  \textbf{Catégorie} & \textbf{Composant} \\
  \midrule
  \texttt{nom} & \texttt{NounDetails} \\
  \texttt{adj}, \texttt{proposs}, \texttt{card}, \texttt{pron} & \texttt{AdjectiveDetails} \\
  \texttt{verb} & \texttt{VerbDetails} \\
  \texttt{conj}, \texttt{part} & \texttt{BaseDetails} \\
  \bottomrule
\end{tabular}
\end{center}

Si le mot possède des phrases d'exemple (\texttt{details.phrases}), un composant \texttt{ExamplePhrases} est affiché en dessous.

\subsubsection{NounDetails.svelte}

\begin{description}
  \item[Rôle] Tableau de déclinaison d'un nom : 6 cas $\times$ 2 nombres (singulier/pluriel).
  \item[Props] \texttt{details} --- l'objet du nom dans \texttt{wordData}.
  \item[Rendu] Pour chaque cas (\ukr{називний}, \ukr{родовий}, \ukr{давальний}, \ukr{знахідний}, \ukr{орудний}, \ukr{місцевий}), une ligne avec la forme singulière et plurielle accentuée.
\end{description}

% ============ PLACEHOLDER ============
% Capture : tableau de déclinaison d'un nom (ex: "будинок")
% montrant les 6 cas avec singulier et pluriel
\begin{figure}[ht]
  \centering
  \fbox{\parbox{0.6\textwidth}{\centering\vspace{2.5cm}\textit{[screenshot-noun-table.png]\\Tableau de déclinaison d'un nom}\vspace{2.5cm}}}
  % \includegraphics[width=0.6\textwidth]{screenshot-noun-table.png}
  \caption{Déclinaison d'un nom --- 6 cas, singulier et pluriel.}
  \label{fig:noun-table}
\end{figure}

\subsubsection{AdjectiveDetails.svelte}

\begin{description}
  \item[Rôle] Tableau de déclinaison : 6 cas $\times$ 4 genres (masculin, féminin, neutre, pluriel). Utilisé pour les adjectifs, pronoms possessifs, cardinaux et pronoms déterminants.
  \item[Props] \texttt{details}.
  \item[Rendu] Un tableau avec les cas en lignes (noms ukrainiens : \ukr{називний}, \ukr{родовий}, etc.) et les genres en colonnes (\ukr{чол. р.}, \ukr{жін. р.}, \ukr{сер. р.}, \ukr{множина}).
\end{description}

\subsubsection{VerbDetails.svelte}

\begin{description}
  \item[Rôle] Tableau de conjugaison complet d'un verbe.
  \item[Props] \texttt{details}.
  \item[Stores] \texttt{wordData} (pour le couple aspectuel).
  \item[Rendu] Structure du tableau~:
  \begin{itemize}
    \item \textbf{Infinitif} (\ukr{Інфінітив}) --- forme de base accentuée.
    \item \textbf{Par temps} (\ukr{Наказовий спосіб}, \ukr{Майбутній час}, \ukr{Теперішній час}, \ukr{Минулий час}) --- pour chaque temps présent dans les données :
    \begin{itemize}
      \item Passé : lignes par genre (\ukr{чол. р.}, \ukr{жін. р.}, \ukr{сер. р.}) avec singulier et pluriel fusionné.
      \item Autres : lignes par personne (\ukr{1 особа}, \ukr{2 особа}, \ukr{3 особа}) avec colonnes \ukr{Однина}/\ukr{Множина}.
    \end{itemize}
    \item \textbf{Forme impersonnelle} (\ukr{Безособова форма}) --- si présente.
    \item \textbf{Couple aspectuel} --- si le verbe a un champ \texttt{coupl}, affiche le verbe de l'aspect opposé.
  \end{itemize}
\end{description}

% ============ PLACEHOLDER ============
% Capture : tableau de conjugaison d'un verbe (ex: "читати")
% montrant infinitif, présent, futur, passé avec toutes les personnes
\begin{figure}[ht]
  \centering
  \fbox{\parbox{0.6\textwidth}{\centering\vspace{3cm}\textit{[screenshot-verb-conj.png]\\Tableau de conjugaison d'un verbe}\vspace{3cm}}}
  % \includegraphics[width=0.6\textwidth]{screenshot-verb-conj.png}
  \caption{Conjugaison complète d'un verbe avec tous les temps disponibles.}
  \label{fig:verb-conj}
\end{figure}

\subsubsection{BaseDetails.svelte}

\begin{description}
  \item[Rôle] Affichage minimal pour les mots invariables (conjonctions, particules).
  \item[Props] \texttt{details}.
  \item[Rendu] Titre «~Forme de base~» suivi de la forme accentuée, calculée avec \texttt{firstPair(details.base)}.
\end{description}

% ----------------------------------------------------------------
\subsection{Phrases}
% ----------------------------------------------------------------

\subsubsection{PhraseList.svelte}

\begin{description}
  \item[Rôle] Page de recherche de phrases (\texttt{/phrases}).
  \item[Props] Aucune.
  \item[Stores] \texttt{phraseData} (lecture).
  \item[État local] \texttt{searchQuery} --- texte saisi dans la barre de recherche.
  \item[Rendu] Barre de recherche flottante (en haut à droite) + liste filtrée. Chaque phrase est rendue via \texttt{HtmlContent} avec sa traduction française et ses références. Le filtrage utilise \texttt{filterPhrases()} (logique AND multi-termes sur les métadonnées \texttt{ref}).
\end{description}

% ============ PLACEHOLDER ============
% Capture : page /phrases avec une recherche active (ex: "L1")
% montrant les phrases filtrées
\begin{figure}[ht]
  \centering
  \fbox{\parbox{0.7\textwidth}{\centering\vspace{2.5cm}\textit{[screenshot-phrase-search.png]\\Recherche filtrée dans les phrases}\vspace{2.5cm}}}
  % \includegraphics[width=0.7\textwidth]{screenshot-phrase-search.png}
  \caption{Recherche de phrases avec filtre actif.}
  \label{fig:phrase-search}
\end{figure}

\subsubsection{ExamplePhrases.svelte}

\begin{description}
  \item[Rôle] Affiche les phrases d'exemple associées à un mot (sous la fiche détaillée).
  \item[Props] \texttt{phrases} --- objet dont les clés sont des identifiants de phrases.
  \item[Stores] \texttt{phraseData}, \texttt{wordData} (lecture).
  \item[État local] \texttt{expandedVerbs} --- objet indiquant quelles phrases ont les formes verbales dépliées.
  \item[Rendu] Pour chaque phrase~: contenu HTML (via \texttt{HtmlContent}), traduction, remarque optionnelle. Si la phrase a un champ \texttt{genereVerbe}, une case à cocher permet de déplier toutes les formes conjuguées (générées par \texttt{generateVerbForms()}).
\end{description}

% ----------------------------------------------------------------
\subsection{Grammaire interactive}
% ----------------------------------------------------------------

\subsubsection{HtmlContent.svelte}

C'est le composant central de l'interactivité. Il reçoit du HTML brut (contenant des spans \texttt{.ukr}) et y applique deux traitements~:

\begin{description}
  \item[Props] \texttt{html} (chaîne HTML), \texttt{disableHover} (booléen, défaut \texttt{false}).
  \item[Stores] \texttt{wordData}, \texttt{accentEnabled}, \texttt{pinnedElement}, \texttt{grammarTableData}.
  \item[Fonctionnement] Un \texttt{\$effect} réactif se déclenche quand \texttt{accentEnabled} ou \texttt{html} change~:
  \begin{enumerate}
    \item \textbf{applyAccents()} --- Parcourt tous les \texttt{.ukr} spans. Si les accents sont activés, utilise \texttt{data-info} pour retrouver la position d'accent dans \texttt{wordData}, puis applique \texttt{highlightLetter()} pour colorer la lettre accentuée en rouge (\textit{crimson}).
    \item \textbf{applyHoverHandlers()} --- \emph{Seulement si} \texttt{disableHover} est \texttt{false}. Attache 3 écouteurs d'événements à chaque span \texttt{.ukr}~:
    \begin{itemize}
      \item \texttt{mouseenter} : colore le mot selon le cas grammatical + affiche la bulle d'info + met à jour \texttt{grammarTableData}.
      \item \texttt{mouseleave} : retire la couleur, cache la bulle, efface \texttt{grammarTableData} (sauf si un élément est épinglé).
      \item \texttt{click} : épingle/désépingle l'élément via \texttt{pinnedElement}.
    \end{itemize}
  \end{enumerate}
\end{description}

\paragraph{Le mécanisme d'accent} L'attribut \texttt{data-info} de chaque span \texttt{.ukr} contient une chaîne séparée par des points-virgules, par exemple~:

\begin{ukrbox}
\texttt{data-info="будинок;nom;cas;nomi;s"}
\end{ukrbox}

Cette chaîne est parsée par \texttt{parseInfo()} pour obtenir : le lemme (\ukr{будинок}), la catégorie (\texttt{nom}), puis les tokens de navigation dans le JSON (\texttt{cas}, \texttt{nomi}, \texttt{s}). La fonction \texttt{getDataFromJson()} utilise ces tokens pour retrouver l'entrée correspondante, d'où on extrait la position d'accent.

% ============ PLACEHOLDER ============
% Capture : info-bulle apparaissant au survol d'un mot ukrainien
% dans le panneau central ou dans une phrase
\begin{figure}[ht]
  \centering
  \fbox{\parbox{0.5\textwidth}{\centering\vspace{2cm}\textit{[screenshot-hover-bubble.png]\\Info-bulle au survol}\vspace{2cm}}}
  % \includegraphics[width=0.5\textwidth]{screenshot-hover-bubble.png}
  \caption{Info-bulle au survol d'un mot ukrainien : lemme accentué + métadonnées.}
  \label{fig:hover-bubble}
\end{figure}

\subsubsection{GrammarSidebar.svelte}

\begin{description}
  \item[Rôle] Sidebar droite affichant le tableau de grammaire complet du mot survolé ou épinglé.
  \item[Props] Aucune.
  \item[Stores] \texttt{grammarTableData}, \texttt{pinnedElement} (lecture).
  \item[Rendu] Visible uniquement si \texttt{grammarTableData} n'est pas \texttt{null}. Affiche un en-tête avec le lemme accentué et les métadonnées (catégorie, genre pour les noms, aspect pour les verbes), puis délègue au composant \texttt{GrammarTable}. Quand un élément est épinglé, la bordure change de couleur (classe CSS \texttt{.pinned}).
\end{description}

% ============ PLACEHOLDER ============
% Captures : sidebar grammaire en mode hover et en mode épinglé
\begin{figure}[ht]
  \centering
  \fbox{\parbox{0.4\textwidth}{\centering\vspace{2.5cm}\textit{[screenshot-grammar-sidebar.png]\\Sidebar grammaire (hover)}\vspace{2.5cm}}}
  % \includegraphics[width=0.4\textwidth]{screenshot-grammar-sidebar.png}
  \caption{Sidebar de grammaire en mode survol.}
  \label{fig:grammar-sidebar}
\end{figure}

\begin{figure}[ht]
  \centering
  \fbox{\parbox{0.4\textwidth}{\centering\vspace{2.5cm}\textit{[screenshot-grammar-pinned.png]\\Sidebar grammaire (épinglée)}\vspace{2.5cm}}}
  % \includegraphics[width=0.4\textwidth]{screenshot-grammar-pinned.png}
  \caption{Sidebar de grammaire épinglée (après clic sur un mot).}
  \label{fig:grammar-pinned}
\end{figure}

\subsubsection{GrammarTable.svelte}

\begin{description}
  \item[Rôle] Génère le tableau HTML de grammaire selon la catégorie du mot.
  \item[Props] \texttt{data} --- objet \texttt{\{word, category, infos, wordData\}}.
  \item[Rendu] Utilise \texttt{\$derived.by()} pour calculer le HTML du tableau. Selon la catégorie~:
  \begin{itemize}
    \item \textbf{Noms} : cas $\times$ nombre, forme courante en gras.
    \item \textbf{Pronoms personnels} : cas simple, forme courante en gras.
    \item \textbf{Adjectifs/pronoms/cardinaux} : cas $\times$ genre, forme courante en gras.
    \item \textbf{Verbes} : infinitif + tous les temps (présent, futur, passé, impératif, impersonnel) + couple aspectuel.
  \end{itemize}
  La fonction interne \texttt{renderCell()} affiche les formes avec accent, séparées par «~/ ~» en cas de variantes multiples.
\end{description}

\subsubsection{UkrSpan.svelte}

\begin{description}
  \item[Rôle] Version composant-propre d'un span ukrainien interactif (alternative à l'approche DOM de \texttt{HtmlContent}).
  \item[Props] \texttt{dataInfo} (chaîne), \texttt{text} (texte à afficher).
  \item[Stores] \texttt{accentEnabled}, \texttt{pinnedElement}, \texttt{grammarTableData}, \texttt{wordData}.
  \item[Rendu] Un \texttt{<span class="ukr">} avec gestion de l'accent, survol, bulle et épinglage, implémentée de manière réactive plutôt que par manipulation DOM.
\end{description}

% ============ PLACEHOLDER ============
% Captures : accents activés vs désactivés
\begin{figure}[ht]
  \centering
  \fbox{\parbox{0.4\textwidth}{\centering\vspace{2cm}\textit{[screenshot-accent-on.png]\\Accents activés}\vspace{2cm}}}
  \hfill
  \fbox{\parbox{0.4\textwidth}{\centering\vspace{2cm}\textit{[screenshot-accent-off.png]\\Accents désactivés}\vspace{2cm}}}
  % \includegraphics[width=0.45\textwidth]{screenshot-accent-on.png}
  % \hfill
  % \includegraphics[width=0.45\textwidth]{screenshot-accent-off.png}
  \caption{Comparaison : accents activés (gauche) vs désactivés (droite).}
  \label{fig:accents}
\end{figure}

\newpage
\section{Stores --- État global}

L'application utilise des \emph{stores} Svelte (\texttt{writable}) pour partager l'état entre composants sans passer par les props. Deux fichiers définissent tous les stores.

\subsection{dataStore.js --- Données du dictionnaire}

\begin{lstlisting}[language=javascript, title=src/lib/stores/dataStore.js]
import { writable } from 'svelte/store';

export const wordData = writable({});
export const phraseData = writable({});
\end{lstlisting}

\begin{description}
  \item[\texttt{wordData}] Contient l'intégralité du dictionnaire, structuré par catégorie grammaticale. Chargé depuis \texttt{static/data.json} au démarrage via \texttt{+layout.server.js} $\rightarrow$ \texttt{+layout.svelte}. Voir l'annexe~\ref{sec:data-schema} pour le schéma complet.

  \item[\texttt{phraseData}] Contient toutes les phrases d'exemple avec traductions et références. Chargé depuis \texttt{static/phrases.json}.
\end{description}

\subsection{uiStore.js --- État de l'interface}

\begin{lstlisting}[language=javascript, title=src/lib/stores/uiStore.js]
import { writable } from 'svelte/store';

export const pinnedElement = writable(null);
export const accentEnabled = writable(true);
export const selectedWord = writable(null);
export const selectedCategory = writable(null);
export const grammarTableData = writable(null);
\end{lstlisting}

\begin{description}
  \item[\texttt{selectedWord}] Le mot actuellement sélectionné dans la sidebar (ex: \texttt{"будинок"}). Mis à jour par \texttt{WordList} au clic.

  \item[\texttt{selectedCategory}] La catégorie du mot sélectionné (ex: \texttt{"nom"}). Toujours mis à jour en même temps que \texttt{selectedWord}.

  \item[\texttt{accentEnabled}] Booléen indiquant si les accents toniques sont affichés. Défaut : \texttt{true}. Basculé par \texttt{AccentCheckbox}.

  \item[\texttt{grammarTableData}] Objet \texttt{\{word, category, infos, wordData\}} transmis à la sidebar de grammaire lors du survol ou du clic sur un span \texttt{.ukr}. Mis à \texttt{null} quand aucun mot n'est survolé (et qu'aucun élément n'est épinglé).

  \item[\texttt{pinnedElement}] Référence DOM de l'élément épinglé (ou \texttt{null}). Quand un élément est épinglé :
  \begin{itemize}
    \item La sidebar de grammaire reste affichée même sans survol.
    \item Les survols d'autres mots ne modifient pas \texttt{grammarTableData}.
    \item Un second clic sur le même élément le désépingle.
  \end{itemize}
\end{description}

\subsection{Diagramme du flux réactif}

\begin{lstlisting}[title=Flux des stores dans l'application]
Utilisateur clique sur un mot (WordList)
  --> selectedWord.set("mot")
  --> selectedCategory.set("nom")
  --> WordDetails reagit ($derived sur les stores)
      --> NounDetails / VerbDetails / ... s'affiche

Utilisateur survole un span .ukr (HtmlContent)
  --> grammarTableData.set({word, category, infos, wordData})
  --> GrammarSidebar reagit ($derived sur grammarTableData)
      --> GrammarTable genere le tableau

Utilisateur clique sur un span .ukr (HtmlContent)
  --> pinnedElement.set(element)
  --> GrammarSidebar ajoute la classe .pinned

Utilisateur coche/decoche "Accents" (AccentCheckbox)
  --> accentEnabled.set(true/false)
  --> HtmlContent reagit ($effect sur $accentEnabled)
      --> applyAccents() re-parcourt tous les .ukr
\end{lstlisting}

\newpage
\section{Fonctions utilitaires}

Les utilitaires sont des fonctions pures, sans effet de bord, testées unitairement avec Vitest. Ils sont situés dans \texttt{src/lib/utils/}.

% ================================================================
\subsection{accent.js --- Gestion des accents toniques}

\subsubsection{\texttt{addAccent(word, accentPosition)}}

Insère un accent aigu combinant (U+0301) après le caractère à la position indiquée (indexation 1-based).

\begin{ukrbox}
\texttt{addAccent("один", 3)} $\rightarrow$ \ukr{оди́н}

\texttt{addAccent("будинок", 3)} $\rightarrow$ \ukr{буди́нок}
\end{ukrbox}

Conventions de position~:
\begin{itemize}
  \item Position > 0 : position normale de l'accent.
  \item Position $-1$ : mot monosyllabique (pas d'accent affiché).
  \item Position $-2$ : accent inconnu.
\end{itemize}

\textbf{Utilisé par :} \texttt{HtmlContent}, \texttt{NounDetails}, \texttt{AdjectiveDetails}, \texttt{VerbDetails}, \texttt{BaseDetails}, \texttt{GrammarTable}, \texttt{dataAccess.js}.

\subsubsection{\texttt{highlightLetter(word, position, classe)}}

Entoure la lettre accentuée dans un \texttt{<span>} avec la classe CSS donnée, en ajoutant l'accent combinant. Position 0-based.

\begin{ukrbox}
\texttt{highlightLetter("один", 2, "accent")}

$\rightarrow$ \texttt{од<span class="accent">и́</span>н}
\end{ukrbox}

La classe \texttt{accent} est stylée en rouge crimson dans \texttt{app.css}.

\textbf{Utilisé par :} \texttt{HtmlContent.applyAccents()}.

% ================================================================
\subsection{parsing.js --- Parsing et transformation de données}

\subsubsection{\texttt{parseInfo(info)}}

Découpe une chaîne \texttt{data-info} par point-virgule.

\begin{ukrbox}
\texttt{parseInfo("будинок;nom;cas;nomi;s")}

$\rightarrow$ \texttt{["будинок", "nom", "cas", "nomi", "s"]}
\end{ukrbox}

\subsubsection{\texttt{toPairs(entry)}}

Normalise les différents formats d'entrée JSON en un tableau de paires \texttt{[[texte, position], ...]}.

Formats reconnus~:
\begin{enumerate}
  \item Paire simple : \texttt{["слово", 3]} $\rightarrow$ \texttt{[["слово", 3]]}
  \item Tableau de paires : \texttt{[["слово", 3], ["варіант", 2]]} $\rightarrow$ inchangé
  \item Ancien format aplati : \texttt{["mot1", 3, "mot2", 3]} $\rightarrow$ \texttt{[["mot1", 3], ["mot2", 3]]}
  \item Ancien format verbe : \texttt{["читаєм", "читаємо", 4]} $\rightarrow$ \texttt{[["читаєм", 4], ["читаємо", 4]]}
\end{enumerate}

\subsubsection{\texttt{firstPair(entry, variantIndex = 0)}}

Retourne la première paire \texttt{[texte, position]} de l'entrée, ou la paire à l'index de variante spécifié. Retourne \texttt{null} si vide.

\subsubsection{\texttt{firstText(entry, variantIndex = 0)}}

Raccourci : extrait le texte de la première paire.

\subsubsection{\texttt{firstAccent(entry, variantIndex = 0)}}

Raccourci : extrait la position d'accent de la première paire.

\subsubsection{\texttt{getVariantIndex(dataInfoTokens)}}

Cherche un token \texttt{var=N} dans les tokens du \texttt{data-info} pour sélectionner une variante spécifique.

\begin{ukrbox}
\texttt{getVariantIndex(["слово", "nom", "cas", "nomi", "s", "var=1"])} $\rightarrow$ \texttt{1}
\end{ukrbox}

\textbf{Utilisé par :} \texttt{HtmlContent}, \texttt{GrammarTable}, \texttt{gramFunc.js}.

% ================================================================
\subsection{dataAccess.js --- Accès aux données du dictionnaire}

\subsubsection{\texttt{getDataFromJson(wordData, category, infos)}}

Navigue dans la structure JSON du dictionnaire pour retrouver l'entrée correspondant à un \texttt{data-info}. La navigation dépend de la catégorie~:

\begin{center}
\begin{tabular}{ll}
  \toprule
  \textbf{Catégorie} & \textbf{Chemin d'accès} \\
  \midrule
  \texttt{nom} & \texttt{wordData.nom[mot][fonction][cas][nombre]} \\
  \texttt{adj/card/proposs/pron} & \texttt{wordData[cat][mot][fonction][cas][genre]} \\
  \texttt{proper} & \texttt{wordData.proper[mot][fonction][cas]} \\
  \texttt{verb} (conjugué) & \texttt{wordData.verb[mot].conj[temps][personne][nombre]} \\
  \texttt{verb} (infinitif) & \texttt{wordData.verb[mot].inf} \\
  \texttt{adv/conj/part} & \texttt{wordData[cat][mot].base} \\
  \texttt{prep} & \texttt{wordData.prep[mot].base} \\
  \bottomrule
\end{tabular}
\end{center}

\textbf{Utilisé par :} \texttt{HtmlContent.applyAccents()}.

\subsubsection{\texttt{getPrincipalForm(wordData, word, category)}}

Retourne la forme canonique accentuée d'un mot (lemme). Utilisée dans l'en-tête de la sidebar de grammaire.

\begin{center}
\begin{tabular}{ll}
  \toprule
  \textbf{Catégorie} & \textbf{Forme canonique} \\
  \midrule
  \texttt{nom} & Nominatif singulier \\
  \texttt{adj/card/proposs/pron} & Nominatif masculin \\
  \texttt{proper} & Nominatif \\
  \texttt{verb} & Infinitif \\
  \bottomrule
\end{tabular}
\end{center}

\textbf{Utilisé par :} \texttt{GrammarSidebar}.

% ================================================================
\subsection{gramFunc.js --- Génération de formes verbales}

\subsubsection{\texttt{generateVerbForms(wordData, verb, tense, fragment1, fragment2)}}

Génère du HTML listant toutes les formes conjuguées d'un verbe pour un temps donné, insérées dans un contexte phrastique (\texttt{fragment1} ... pronom ... forme ... \texttt{fragment2}).

Pour le passé (\texttt{pass}), itère par genre (m, f, n) avec les pronoms \ukr{він}/\ukr{вона}/\ukr{воно}. Pour les autres temps, itère par personne (1p, 2p, 3p) avec les pronoms \ukr{я}/\ukr{ти}/\ukr{він}$\cdots$

Chaque forme est un span \texttt{.ukr} cliquable avec son \texttt{data-info}, ce qui permet les interactions de hover/pin.

\textbf{Utilisé par :} \texttt{ExamplePhrases}.

% ================================================================
\subsection{colors.js --- Couleurs par cas grammatical}

Exporte un objet \texttt{classesToColors} associant chaque cas à une couleur RGB~:

\begin{center}
\begin{tabular}{lll}
  \toprule
  \textbf{Cas} & \textbf{Clé} & \textbf{Couleur} \\
  \midrule
  Nominatif & \texttt{nomi} & inherit (noir) \\
  Génitif & \texttt{gen} & \textcolor[rgb]{0.07,0.52,0.18}{vert} \\
  Accusatif & \texttt{acc} & \textcolor[rgb]{0.22,0.02,0.97}{bleu} \\
  Locatif & \texttt{loc} & \textcolor[rgb]{0.52,0.36,0.07}{brun} \\
  Instrumental/Datif & \texttt{ins}/\texttt{dat} & \textcolor[rgb]{0.91,0.16,0.61}{rose} \\
  Vocatif & \texttt{voc} & \textcolor[rgb]{0.52,0.07,0.07}{rouge foncé} \\
  \bottomrule
\end{tabular}
\end{center}

\textbf{Utilisé par :} \texttt{HtmlContent.applyHoverHandlers()}.

% ================================================================
\subsection{i18n.js --- Labels multilingues}

Fournit des traductions courtes pour les catégories grammaticales, temps et nombres~:

\begin{lstlisting}[language=javascript]
labelCategory("verb")  // -> "verbe"
labelCategory("nom")   // -> "nom"
labelTense("pres")     // -> "pres."
labelNumber("s")       // -> "sg"
\end{lstlisting}

\textbf{Utilisé par :} \texttt{HtmlContent.buildBubbleHTML()}, \texttt{GrammarSidebar}, \texttt{GrammarTable}.

% ================================================================
\subsection{foldState.js --- État de pliage de la sidebar}

\subsubsection{\texttt{hasAnyExpanded(groupedData, categoryOpen, letterOpen)}}

Retourne \texttt{true} si au moins une catégorie est ouverte ou un groupe de lettres est déplié. Fonction pure utilisée par \texttt{WordList.svelte} pour piloter le label du bouton global «~Tout déplier / Tout replier~».

\begin{lstlisting}[language=javascript]
hasAnyExpanded(groupedData, {nom: true, adj: false}, {})
// -> true (la categorie "nom" est ouverte)

hasAnyExpanded(groupedData, {nom: false}, {"nom:A": true})
// -> true (le groupe de lettres "A" est ouvert)

hasAnyExpanded(groupedData, {nom: false}, {"nom:A": false})
// -> false (tout est ferme)
\end{lstlisting}

\textbf{Utilisé par :} \texttt{WordList.svelte} (valeur dérivée \texttt{anyExpanded}).

% ================================================================
\subsection{phrases.js --- Filtrage de phrases}

\subsubsection{\texttt{filterPhrases(phrasesData, searchQuery)}}

Filtre les phrases par une recherche multi-termes (logique AND). Chaque terme doit apparaître dans les métadonnées \texttt{ref} de la phrase (clés et valeurs concaténées, insensible à la casse).

\begin{ukrbox}
\texttt{filterPhrases(data, "L1 Cours5")} $\rightarrow$ uniquement les phrases dont \texttt{ref} contient à la fois «~L1~» et «~Cours5~».
\end{ukrbox}

\textbf{Utilisé par :} \texttt{PhraseList}.

\newpage
\section{Pipeline de données}

Les données affichées par l'application proviennent de fichiers JSON statiques, générés par des scripts Python à partir de sources externes (dictionnaire Goroh en ligne) et de phrases annotées manuellement.

\subsection{Vue d'ensemble}

\begin{lstlisting}[title=Pipeline de generation des donnees]
phrases_a_traiter.json    (phrases annotees avec spans .ukr)
        |
        v
build_entries_2_(nooj_opt).py   (scraping + parsing morphologique)
        |
        v
    out.json              (nouvelles entrees generees)
        |
        v  (fusion manuelle)
  static/data.json        (dictionnaire complet, utilise par l'app)
  static/phrases.json     (phrases d'exemple)
\end{lstlisting}

\subsection{Fichiers du pipeline}

Tous situés dans \texttt{outil\_python/}~:

\begin{description}
  \item[\texttt{build\_entries\_2\_(nooj\_opt).py}] Script principal. Pour chaque mot ukrainien trouvé dans les phrases à traiter~:
  \begin{enumerate}
    \item Vérifie s'il existe déjà dans \texttt{data.json} (dédoublement).
    \item Cherche le lemme et la catégorie grammaticale.
    \item Scrape le dictionnaire Goroh (\texttt{goroh.ua}) pour obtenir les tables de déclinaison/conjugaison.
    \item Parse les tables HTML en structure JSON normalisée.
    \item Génère le champ \texttt{base\_html} (span canonique avec \texttt{data-info}).
    \item Écrit le résultat dans \texttt{out.json}.
  \end{enumerate}

  \item[\texttt{phrases\_a\_traiter.json}] Fichier d'entrée : liste de phrases avec leurs spans \texttt{.ukr} annotés (traduction, références, remarques).

  \item[\texttt{out.json}] Fichier de sortie : nouvelles entrées prêtes à être fusionnées dans \texttt{static/data.json}.

  \item[\texttt{entries\_report.html}] Rapport HTML des entrées traitées (pour vérification manuelle).
\end{description}

\subsection{Format des données générées}

Chaque entrée dans \texttt{out.json} suit le format attendu par l'application (voir annexe~\ref{sec:data-schema}). Points importants~:

\begin{itemize}
  \item \textbf{Ordre des champs} : \texttt{cas} (ou \texttt{inf}/\texttt{conj} pour les verbes), puis \texttt{genre} (noms), \texttt{asp} (verbes), \texttt{nooj} (optionnel), \texttt{base\_html}, \texttt{phrases}.
  \item \textbf{Collisions de lemmes} : si deux mots ont le même lemme mais des catégories différentes, le suffixe \texttt{\_\_pos} est ajouté (ex: \texttt{один\_\_card} vs \texttt{один\_\_adj}).
  \item \textbf{Format canonique de \texttt{base\_html}} :
\end{itemize}

\begin{lstlisting}[language=html]
<span class="ukr" data-info="LEMME;CATEGORIE;...">TEXTE</span>
\end{lstlisting}

\subsection{Mise à jour des données}

Le processus de mise à jour est semi-automatique~:

\begin{enumerate}
  \item Ajouter les nouvelles phrases dans \texttt{phrases\_a\_traiter.json}.
  \item Exécuter le script Python : \texttt{python build\_entries\_2\_(nooj\_opt).py}.
  \item Vérifier \texttt{entries\_report.html} pour les erreurs.
  \item Fusionner manuellement \texttt{out.json} dans \texttt{static/data.json}.
  \item Ajouter les phrases dans \texttt{static/phrases.json}.
  \item Rebuilder l'application : \texttt{npm run build}.
\end{enumerate}


\appendix
\newpage
\section{Annexe : SvelteKit pour les non-initiés}
\label{sec:sveltekit}

Cette annexe explique les concepts SvelteKit utilisés dans l'application, à destination de lecteurs familiers avec HTML/CSS/JS mais pas avec les frameworks modernes.

\subsection{Qu'est-ce qu'un composant Svelte~?}

Un fichier \texttt{.svelte} est un composant autonome qui regroupe HTML, JavaScript et CSS en un seul fichier~:

\begin{lstlisting}[language=svelte, title=Exemple : un composant minimal]
<script>
  let count = $state(0);
</script>

<button onclick={() => count++}>
  Compteur : {count}
</button>

<style>
  button { font-size: 1.2rem; }
</style>
\end{lstlisting}

\begin{itemize}
  \item \texttt{<script>} : logique JavaScript.
  \item Corps HTML : le template, avec des expressions \texttt{\{...\}} pour les valeurs dynamiques.
  \item \texttt{<style>} : CSS \emph{scopé} --- il ne s'applique qu'à ce composant, pas aux autres.
\end{itemize}

\subsection{Réactivité Svelte~5}

Svelte~5 utilise des \emph{runes} (fonctions spéciales préfixées par \texttt{\$}) pour la réactivité~:

\subsubsection{\texttt{\$state} --- État local réactif}

Déclare une variable dont les modifications sont automatiquement reflétées dans le DOM.

\begin{lstlisting}[language=javascript]
let count = $state(0);   // reactif
count++;                  // le DOM se met a jour automatiquement
\end{lstlisting}

\textbf{Dans l'application :} \texttt{HtmlContent} utilise \texttt{let container = \$state(null)} pour stocker la référence DOM.

\subsubsection{\texttt{\$derived} --- Valeur dérivée}

Calcule une valeur qui se met à jour automatiquement quand ses dépendances changent.

\begin{lstlisting}[language=javascript]
let count = $state(0);
let doubled = $derived(count * 2);  // toujours = count * 2
\end{lstlisting}

\textbf{Dans l'application :} \texttt{WordDetails} utilise \texttt{\$derived} pour calculer les détails du mot sélectionné à partir des stores.

La variante \texttt{\$derived.by(() => ...)} permet des calculs plus complexes (multi-lignes).

\subsubsection{\texttt{\$effect} --- Effet de bord réactif}

Exécute du code quand ses dépendances réactives changent.

\begin{lstlisting}[language=javascript]
$effect(() => {
  const enabled = $accentEnabled;  // dependance trackee
  if (container) applyAccents(container);
});
\end{lstlisting}

S'exécute :
\begin{enumerate}
  \item Au montage initial du composant.
  \item À chaque fois qu'une dépendance trackée change.
\end{enumerate}

\textbf{Dans l'application :} \texttt{HtmlContent} utilise \texttt{\$effect} pour ré-appliquer les accents et les handlers de hover quand l'état change.

\subsubsection{\texttt{\$props} --- Props du composant}

Déclare les propriétés reçues du composant parent.

\begin{lstlisting}[language=javascript]
// Dans le composant enfant :
let { html, disableHover = false } = $props();

// Dans le composant parent :
<HtmlContent html={wordInfo.base_html} disableHover={true} />
\end{lstlisting}

\subsection{Stores Svelte}

Un \emph{store} est un conteneur de valeur partagé entre composants, sans avoir besoin de passer des props à travers l'arbre.

\begin{lstlisting}[language=javascript, title=Creer un store]
import { writable } from 'svelte/store';
export const count = writable(0);
\end{lstlisting}

\begin{lstlisting}[language=javascript, title=Utiliser un store]
// Ecriture :
count.set(42);

// Lecture dans un composant (syntaxe $) :
<p>Valeur : {$count}</p>

// Lecture hors composant :
import { get } from 'svelte/store';
const val = get(count);
\end{lstlisting}

La syntaxe \texttt{\$count} (préfixe \texttt{\$}) est un raccourci Svelte qui s'abonne automatiquement au store et réagit aux changements. Hors d'un composant (dans une fonction utilitaire), il faut utiliser \texttt{get(store)}.

\textbf{Dans l'application :} Les stores sont définis dans \texttt{src/lib/stores/} et utilisés partout --- \texttt{\$wordData} dans les composants, \texttt{get(wordData)} dans \texttt{HtmlContent} (fonctions impératives).

\subsection{Routage basé sur le système de fichiers}

SvelteKit associe automatiquement les URLs aux fichiers dans \texttt{src/routes/}~:

\begin{center}
\begin{tabular}{ll}
  \toprule
  \textbf{Fichier} & \textbf{URL} \\
  \midrule
  \texttt{src/routes/+page.svelte} & \texttt{/} \\
  \texttt{src/routes/phrases/+page.svelte} & \texttt{/phrases} \\
  \bottomrule
\end{tabular}
\end{center}

Fichiers spéciaux dans chaque dossier de route~:
\begin{description}
  \item[\texttt{+page.svelte}] Le contenu de la page.
  \item[\texttt{+layout.svelte}] Template partagé par toutes les pages du dossier (et sous-dossiers). Affiche \texttt{\{@render children()\}} là où le contenu de la page doit apparaître.
  \item[\texttt{+layout.server.js}] Code exécuté côté serveur (ou au build) pour charger des données. La fonction \texttt{load()} retourne un objet accessible dans le layout via \texttt{\$props().data}.
\end{description}

\subsection{Adapter-static et prérendu}

Par défaut, SvelteKit génère une application serveur. L'option \texttt{adapter-static} transforme toutes les pages en HTML statique au moment du build~:

\begin{lstlisting}[language=javascript, title=svelte.config.js]
import adapter from '@sveltejs/adapter-static';

export default {
  kit: {
    adapter: adapter({
      pages: 'build',
      assets: 'build',
      strict: true
    })
  }
};
\end{lstlisting}

Avec \texttt{export const prerender = true} dans \texttt{+layout.server.js}, toutes les pages sont générées au build. Le résultat dans \texttt{build/} peut être déployé sur n'importe quel serveur statique (Apache, Nginx, GitHub Pages, etc.).

\subsection{Comparaison avant/après migration}

\begin{center}
\begin{tabular}{p{0.45\textwidth}p{0.45\textwidth}}
  \toprule
  \textbf{Vanilla JS (avant)} & \textbf{SvelteKit (après)} \\
  \midrule
  Scripts dans \texttt{scripts/*.js} avec manipulation DOM manuelle & Composants \texttt{.svelte} déclaratifs \\
  \midrule
  État global via variables globales & Stores Svelte (\texttt{writable}) \\
  \midrule
  \texttt{innerHTML} pour les tableaux & Templates Svelte avec \texttt{\{\#each\}}, \texttt{\{\#if\}} \\
  \midrule
  \texttt{document.querySelector} pour cibler les zones & Binding avec \texttt{bind:this} et refs \\
  \midrule
  Pages HTML séparées (\texttt{index.html}, \texttt{affiche\_phrase.html}) & Routes SvelteKit (\texttt{/}, \texttt{/phrases}) \\
  \midrule
  Chargement JSON via \texttt{fetch} au runtime & Prérendu : JSON intégré au build \\
  \midrule
  CSS global unique (\texttt{styles.css}) & CSS scopé par composant + \texttt{app.css} global \\
  \bottomrule
\end{tabular}
\end{center}

\newpage
\section{Annexe : Schéma des données JSON}
\label{sec:data-schema}

\subsection{data.json --- Dictionnaire}

Le fichier \texttt{static/data.json} contient l'intégralité du dictionnaire, organisé par catégorie grammaticale au premier niveau.

\subsubsection{Structure générale}

\begin{lstlisting}[title=Structure de premier niveau]
{
  "prep":    { ... },    // Prepositions
  "nom":     { ... },    // Noms
  "adj":     { ... },    // Adjectifs
  "card":    { ... },    // Cardinaux
  "proposs": { ... },    // Pronoms possessifs
  "pron":    { ... },    // Pronoms (determinants)
  "proper":  { ... },    // Pronoms personnels
  "verb":    { ... },    // Verbes
  "adv":     { ... },    // Adverbes
  "conj":    { ... },    // Conjonctions
  "part":    { ... }     // Particules
}
\end{lstlisting}

\subsubsection{Format des formes fléchies}

Chaque forme fléchie est stockée sous l'un de ces formats~:

\begin{enumerate}
  \item \textbf{Paire simple} : \texttt{["текст", position]}

  Exemple : \texttt{["будинку", 3]} signifie \ukr{буди́нку} (accent sur la 3\ieme{} lettre).

  \item \textbf{Variantes multiples} : \texttt{[["варіант1", pos1], ["варіант2", pos2]]}

  Pour les mots ayant plusieurs formes acceptées. La sélection se fait via \texttt{var=N} dans \texttt{data-info}.

  \item \textbf{Ancien format (legacy)} : \texttt{["форма1", "форма2", position]}

  Plusieurs formes partageant la même position d'accent. Converti automatiquement par \texttt{toPairs()}.
\end{enumerate}

Conventions de position~:
\begin{itemize}
  \item $> 0$ : position de l'accent (1-indexed).
  \item $-1$ : mot monosyllabique (pas d'accent visible).
  \item $-2$ : accent inconnu.
\end{itemize}

\subsubsection{Noms (\texttt{nom})}

\begin{lstlisting}[language=javascript]
"nom": {
  "LEMME": {
    "cas": {
      "nomi": { "s": ["forme_s", pos], "pl": ["forme_pl", pos] },
      "gen":  { "s": [...], "pl": [...] },
      "dat":  { "s": [...], "pl": [...] },
      "acc":  { "s": [...], "pl": [...] },
      "ins":  { "s": [...], "pl": [...] },
      "loc":  { "s": [...], "pl": [...] }
    },
    "genre": "m" | "f" | "n",
    "base_html": "<span class=\"ukr\" data-info=\"LEMME;nom;cas;nomi;s\">...</span>",
    "phrases": { "phraseKey1": true, ... }   // optionnel
  }
}
\end{lstlisting}

Les 6 cas ukrainiens~:
\begin{center}
\begin{tabular}{lll}
  \toprule
  \textbf{Clé} & \textbf{Cas} & \textbf{Ukrainien} \\
  \midrule
  \texttt{nomi} & Nominatif & \ukr{називний} \\
  \texttt{gen} & Génitif & \ukr{родовий} \\
  \texttt{dat} & Datif & \ukr{давальний} \\
  \texttt{acc} & Accusatif & \ukr{знахідний} \\
  \texttt{ins} & Instrumental & \ukr{орудний} \\
  \texttt{loc} & Locatif & \ukr{місцевий} \\
  \bottomrule
\end{tabular}
\end{center}

\subsubsection{Adjectifs, pronoms possessifs, cardinaux, pronoms (\texttt{adj}, \texttt{proposs}, \texttt{card}, \texttt{pron})}

Même structure que les noms, mais avec 4 genres au lieu de 2 nombres~:

\begin{lstlisting}[language=javascript]
"adj": {
  "LEMME": {
    "cas": {
      "nomi": { "m": [...], "f": [...], "n": [...], "pl": [...] },
      "gen":  { "m": [...], "f": [...], "n": [...], "pl": [...] },
      // ... autres cas
    },
    "base_html": "...",
    "phrases": { ... }
  }
}
\end{lstlisting}

\subsubsection{Pronoms personnels (\texttt{proper})}

Pas de genre dans la déclinaison --- seulement les cas~:

\begin{lstlisting}[language=javascript]
"proper": {
  "LEMME": {
    "cas": {
      "nomi": ["forme", pos],
      "gen":  ["forme", pos],
      // ...
    },
    "base_html": "..."
  }
}
\end{lstlisting}

\subsubsection{Verbes (\texttt{verb})}

\begin{lstlisting}[language=javascript]
"verb": {
  "LEMME": {
    "inf": ["infinitif", pos],
    "conj": {
      "pres": {           // present (verbes imperfectifs)
        "1p": { "s": [...], "pl": [...] },
        "2p": { "s": [...], "pl": [...] },
        "3p": { "s": [...], "pl": [...] }
      },
      "fut": { ... },     // futur (meme structure que pres)
      "pass": {            // passe (structure par genre)
        "m":  { "s": [...], "pl": [...] },
        "f":  { "s": [...], "pl": [...] },
        "n":  { "s": [...], "pl": [...] }
      },
      "imp": {             // imperatif (seulement 1p et 2p)
        "1p": { "s": [...], "pl": [...] },
        "2p": { "s": [...], "pl": [...] }
      },
      "impers": [...]      // forme impersonnelle (optionnel)
    },
    "asp": "perfectif" | "imperfectif",
    "coupl": "LEMME_ASPECT_OPPOSE",  // optionnel
    "base_html": "...",
    "phrases": { ... }
  }
}
\end{lstlisting}

\subsubsection{Mots invariables (\texttt{adv}, \texttt{conj}, \texttt{part}, \texttt{prep})}

\begin{lstlisting}[language=javascript]
"adv": {
  "LEMME": {
    "base": ["forme", pos],
    "base_html": "..."
  }
}
\end{lstlisting}

% ================================================================
\subsection{phrases.json --- Phrases d'exemple}

\begin{lstlisting}[language=javascript]
{
  "CLE_PHRASE": {
    "phrase_html": "<span class=\"ukr\" data-info=\"...\">mot</span> ...",
    "traduction": "Traduction francaise",
    "ref": {
      "Auteur": "L1,Cours5,v1.0"
    },
    "remarque": "Note optionnelle",
    "genereVerbe": {             // optionnel
      "verbe": "LEMME_VERBE",
      "temps": "pres",
      "frag1": "texte avant",
      "frag2": "texte apres"
    }
  }
}
\end{lstlisting}

\begin{description}
  \item[\texttt{phrase\_html}] Le contenu HTML de la phrase, avec des spans \texttt{.ukr} interactifs.
  \item[\texttt{traduction}] Traduction française de la phrase.
  \item[\texttt{ref}] Métadonnées de référence (auteur, cours, source). Utilisées pour le filtrage.
  \item[\texttt{remarque}] Note optionnelle affichée sous la phrase.
  \item[\texttt{genereVerbe}] Si présent, permet de générer toutes les formes conjuguées du verbe pour le temps indiqué, insérées dans le contexte \texttt{frag1 ... pronom ... forme ... frag2}.
\end{description}


\end{document}
